\section{Protocol}

\label{sec:protocol}
%% ============================================================================

Quasar 3.0 inherits the Quasar consensus family of staged probabilistic-DAG
protocols. Two engine variants share the family:

\begin{itemize}[leftmargin=2em,nosep]
\item \textbf{Nova} (linear chain): Photon $\to$ Wave $\to$ Focus $\to$ Ray.
\item \textbf{Nebula} (DAG): Flare $\to$ Horizon $\to$ Field, with the
      same QuasarCert at the finalization step.
\end{itemize}

Both variants finalize a block by producing the same QuasarCert
artifact (Section~\ref{sec:quasarcert}). What differs is the upstream
voting structure---linear in Nova, DAG in Nebula. We describe Nova here
and reference~\cite{lux-nova,lux-nebula} for the DAG case.

\subsection{Phase 1: Photon (Proposal Emission)}

The leader for height $h$ (round-robin over the validator set, rotated
each epoch) constructs a block $B_h$ containing:
\begin{itemize}[leftmargin=2em,nosep]
\item The parent block hash $h_{\text{parent}} = H(B_{h-1})$.
\item A batch of transactions from the mempool.
\item The leader's BLS signature over $B_h$.
\item A timestamp $\tau_h$ and the seven chain roots $\boldsymbol{\rho}$
      (Definition~\ref{def:roots}) observed at $\tau_h$.
\end{itemize}
The leader broadcasts $\langle \textsc{Propose}, B_h, \boldsymbol{\rho}
\rangle$ to all validators. Photon emission is fire-and-forget: the leader
does not wait for acknowledgments.

\subsection{Phase 2: Wave (Fast Probabilistic Consensus)}

Upon receiving a valid proposal, each validator $v_i$ enters the Wave
phase. Wave implements Fast Probabilistic Consensus (FPC) with
phase-dependent thresholds.

\begin{definition}[FPC Threshold Selection]
\label{def:fpc}
For round $r$ with sample size $k$, the threshold $\alpha_r$ is:
\[
\alpha_r = \left\lceil \theta_r \cdot k \right\rceil, \qquad
\theta_r = \theta_{\min} + (\theta_{\max} - \theta_{\min}) \cdot
\frac{\mathrm{SHA256}(\mathit{seed} \,\|\, r) \bmod 2^{16}}{2^{16}}.
\]
Default: $\theta_{\min} = 0.5$, $\theta_{\max} = 0.8$.
\end{definition}

\begin{algorithm}
\caption{Wave: FPC voting round}
\label{alg:wave}
\begin{algorithmic}[1]
\Require Validator $v_i$, block $B_h$, round $r$, sample size $k$
\State $S \gets \text{SamplePeers}(k)$
\State $\alpha_r \gets \lceil \text{FPCThreshold}(r, k) \rceil$
\State $\mathit{votes} \gets 0$
\For{$v_j \in S$}
    \State $\mathit{votes} \gets \mathit{votes} +
      \text{QueryPreference}(v_j, B_h)$
\EndFor
\If{$\mathit{votes} \geq \alpha_r$}
    \State $\mathit{pref}_i \gets \textsc{Accept}$
\Else
    \State $\mathit{pref}_i \gets \textsc{Reject}$
\EndIf
\If{$\mathit{pref}_i = \mathit{prevPref}_i$}
    \State $\mathit{confidence}_i \gets \mathit{confidence}_i + 1$
\Else
    \State $\mathit{confidence}_i \gets 0$
\EndIf
\If{$\mathit{confidence}_i \geq \beta$}
    \State \textbf{enter} Focus with $\mathit{pref}_i$
\EndIf
\end{algorithmic}
\end{algorithm}

\subsection{Phase 3: Focus (Certificate Subject Lock)}

Once $\mathit{pref}_i = \textsc{Accept}$ has stabilized for $\beta$
consecutive rounds, $v_i$ enters Focus and computes the certificate
subject of Definition~\ref{def:subject}, binding the seven chain roots
$\boldsymbol{\rho}$, mode $\mu$, parent block hash, parent state root,
parent execution root, gas limit, and base fee. Focus is the boundary
between probabilistic vote propagation (Wave) and cryptographic
finalization (Ray). At this point all three signing lanes share a single
32-byte digest.

\subsection{Phase 4: Ray (Three-Lane Signing)}

In Ray, $v_i$ produces three signatures over the same
$\mathsf{cert\_subject}$ in parallel:

\begin{enumerate}[leftmargin=2em]
\item \textbf{BLS}: $\sigma_i^{\mathrm{BLS}} = \mathit{sk}_i^{\mathrm{BLS}}
      \cdot H_{\mathbb{G}_1}(\mathsf{cert\_subject})$.
      96-byte (uncompressed) or 48-byte (compressed) $\mathbb{G}_1$
      signature.
\item \textbf{Ringtail}: contribute share $\sigma_i^{\mathrm{RT}}$ to
      Q-Chain's two-round threshold protocol on
      $\mathsf{cert\_subject}$ using $\mathit{sk}_i^{\mathrm{RT}}$. The
      Q-Chain combiner produces $\sigma^{\mathrm{RT}}$.
\item \textbf{ML-DSA-65}: $\sigma_i^{\mathrm{ML}} =
      \mathsf{MLDSA.Sign}(\mathit{sk}_i^{\mathrm{ML}},
      \mathsf{cert\_subject})$. 3{,}309-byte signature, transmitted only
      to Z-Chain (not gossiped network-wide).
\end{enumerate}

The validator broadcasts $\langle \textsc{QuasarVote}, i,
\sigma_i^{\mathrm{BLS}}, \sigma_i^{\mathrm{RT}} \rangle$ on the consensus
plane and submits $\sigma_i^{\mathrm{ML}}$ to Z-Chain. Z-Chain proves the
batch (Section~\ref{sec:lane3}) and emits the Groth16 artifact
$(\pi_{\mathrm{ZK}}, h_x)$ where $\pi_{\mathrm{ZK}}$ is the 192-byte
SNARK and $h_x$ is the 32-byte public-inputs hash.

\subsection{Finalization}

A block proposer that observes (a) an aggregate
$\sigma^{\mathrm{BLS}}_{\mathrm{agg}}$ over $\geq 2f+1$ BLS shares,
(b) a valid Ringtail $\sigma^{\mathrm{RT}}$ from Q-Chain, and (c) a valid
$(\pi_{\mathrm{ZK}}, h_x)$ from Z-Chain, constructs the QuasarCert
$\mathcal{C}_h = (\sigma^{\mathrm{BLS}}_{\mathrm{agg}},
\sigma^{\mathrm{RT}}, \pi_{\mathrm{ZK}}, h_x)$ and broadcasts
$\langle \textsc{Finalize}, B_h, \mathcal{C}_h \rangle$. Any validator
verifies $\mathcal{C}_h$ by checking the three lanes against
$\mathsf{cert\_subject}$ as in Definition~\ref{def:verify-qcert}.

\subsection{Epoch Rotation}

Quasar 3.0 groups blocks into epochs of fixed length $E$. At epoch
boundaries:
\begin{itemize}[leftmargin=2em,nosep]
\item Validator sets may change (additions, removals, stake changes);
      $P_\text{val}$ updates.
\item Q-Chain runs a new Ringtail DKG; $Q_\text{cer}$ updates and
      the previous threshold key is retired.
\item Z-Chain may rotate the verifying key (for example, when adding
      validators changes the circuit's public-input arity);
      $Z_\text{vk}$ updates.
\item Epoch certificates summarize the finalized chain state.
\item Epochs older than 6 are pruned from active memory; only epoch
      certificates are retained.
\end{itemize}
The 6-epoch retention window bounds active memory to $O(6 \cdot E \cdot
n)$ regardless of total chain length.


%% ============================================================================
